% GENERATED FILE. Edit canonical lessons, then run npm run generate:books.
% canonical-source-hash: fnv1a64:e2403204b93021a8
% canonical-lessons: GE-C303-tod, GE-C303-beerdigen, GE-C303-beerdigung, GE-C303-sarg, GE-C303-friedhof

\chapter{Die Beerdigung}
\label{ch:ge-die-beerdigung}
\hlchaptermodality{303}

\begin{chapteropening}
\textbf{By the end of this chapter:} \emph{I can talk about a death: death, to bury, a funeral, a coffin and a cemetery.}

Complete the last lesson of chapter 303: I can talk about a death: death, to bury, a funeral, a coffin and a cemetery.
\end{chapteropening}

\section[der Tod]{der Tod — the death}

\label{lesson:GE-C303-tod}



\begin{quote}
Before the new one: say the German for the anniversary, then the German for the farewell.
\end{quote}

\subsection*{You'll want to know: der Tod}
\textbf{der Tod} — \textquotedblleft{}the death\textquotedblright{}.

\begin{quote}
Nach dem Tod seiner Frau lebte er allein.
\end{quote}

\begin{quote}
\emph{After the death of his wife he lived alone.}
\end{quote}

The plural is rare. The adjective is \emph{tot}, already known.

\begin{cousinweb}
\emph{Tod} is inherited from Old High German \emph{tōd}, from Proto-Germanic \emph{dauþuz} \textquotedblleft{}death\textquotedblright{} (EtymDB), the same source as English \emph{death}.
\end{cousinweb}

\begin{grammarlens}[title={What you've built}]
One more word for telling, explaining and describing at B1.
\end{grammarlens}

\subsection*{Your turn}
\begin{itemize}
\raggedright
  \item \emph{Say it:} \emph{der Tod}
  \item \emph{Say it:} \emph{der Tod}, once more
  \item \emph{Say it:} a sentence of your own with \emph{der Tod}
  \item \emph{Recall:} say the German for the anniversary, then the German for the farewell, then say \emph{der Tod} again
\end{itemize}

\begin{tcolorbox}[breakable,colback=teal!4,colframe=teal!35!black,title={Before you move on}]
What does der Tod mean? (\textquotedblleft{}The death\textquotedblright{}.) Say it once more.
\end{tcolorbox}
\section[beerdigen]{beerdigen — to bury}

\label{lesson:GE-C303-beerdigen}



\begin{quote}
Before the new one: say the German for the farewell, then the German for the death.
\end{quote}

\subsection*{You'll want to know: beerdigen}
\textbf{beerdigen} — \textquotedblleft{}to bury\textquotedblright{}.

\begin{quote}
Sie wurde in ihrem Heimatdorf beerdigt.
\end{quote}

\begin{quote}
\emph{She was buried in her home village.}
\end{quote}

Built on \emph{die Erde}. \emph{Begraben} means the same.

\begin{cousinweb}
\emph{beerdigen} is built on \emph{Erde} \textquotedblleft{}earth\textquotedblright{}, an inherited Germanic word with English \emph{earth} and Gothic \emph{airþa} among its cousins (EtymDB): to put into the earth.
\end{cousinweb}

\begin{grammarlens}[title={What you've built}]
One more word for telling, explaining and describing at B1.
\end{grammarlens}

\subsection*{Your turn}
\begin{itemize}
\raggedright
  \item \emph{Say it:} \emph{beerdigen}
  \item \emph{Say it:} \emph{beerdigen}, once more
  \item \emph{Say it:} a sentence of your own with \emph{beerdigen}
  \item \emph{Recall:} say the German for the farewell, then the German for the death, then say \emph{beerdigen} again
\end{itemize}

\begin{tcolorbox}[breakable,colback=teal!4,colframe=teal!35!black,title={Before you move on}]
What does beerdigen mean? (\textquotedblleft{}To bury\textquotedblright{}.) Say it once more.
\end{tcolorbox}
\section[die Beerdigung]{die Beerdigung — the funeral}

\label{lesson:GE-C303-beerdigung}



\begin{quote}
Before the new one: say the German for the death, then the German for to bury.
\end{quote}

\subsection*{You'll want to know: die Beerdigung}
\textbf{die Beerdigung} — \textquotedblleft{}the funeral\textquotedblright{}.

\begin{quote}
Zur Beerdigung kamen viele Leute.
\end{quote}

\begin{quote}
\emph{Many people came to the funeral.}
\end{quote}

Plural: \emph{die Beerdigungen}. In Austria also \emph{das Begräbnis}.

\begin{cousinweb}
\emph{Beerdigung} is built on \emph{beerdigen} \textquotedblleft{}to bury\textquotedblright{}, from \emph{Erde} \textquotedblleft{}earth\textquotedblright{}, an inherited Germanic word (EtymDB).
\end{cousinweb}

\begin{grammarlens}[title={What you've built}]
One more word for telling, explaining and describing at B1.
\end{grammarlens}

\subsection*{Your turn}
\begin{itemize}
\raggedright
  \item \emph{Say it:} \emph{die Beerdigung}
  \item \emph{Say it:} \emph{die Beerdigung}, once more
  \item \emph{Say it:} a sentence of your own with \emph{die Beerdigung}
  \item \emph{Recall:} say the German for the death, then the German for to bury, then say \emph{die Beerdigung} again
\end{itemize}

\begin{tcolorbox}[breakable,colback=teal!4,colframe=teal!35!black,title={Before you move on}]
What does die Beerdigung mean? (\textquotedblleft{}The funeral\textquotedblright{}.) Say it once more.
\end{tcolorbox}
\section[der Sarg]{der Sarg — the coffin}

\label{lesson:GE-C303-sarg}



\begin{quote}
Before the new one: say the German for to bury, then the German for the funeral.
\end{quote}

\subsection*{You'll want to know: der Sarg}
\textbf{der Sarg} — \textquotedblleft{}the coffin\textquotedblright{}.

\begin{quote}
Der Sarg war aus hellem Holz.
\end{quote}

\begin{quote}
\emph{The coffin was made of light wood.}
\end{quote}

Plural: \emph{die Särge}.

\begin{cousinweb}
\emph{Sarg} is inherited from Old High German \emph{saruh} (EtymDB), which is usually traced to Late Latin \emph{sarcophagus} \textquotedblleft{}a stone coffin\textquotedblright{}.
\end{cousinweb}

\begin{grammarlens}[title={What you've built}]
One more word for telling, explaining and describing at B1.
\end{grammarlens}

\subsection*{Your turn}
\begin{itemize}
\raggedright
  \item \emph{Say it:} \emph{der Sarg}
  \item \emph{Say it:} \emph{der Sarg}, once more
  \item \emph{Say it:} a sentence of your own with \emph{der Sarg}
  \item \emph{Recall:} say the German for to bury, then the German for the funeral, then say \emph{der Sarg} again
\end{itemize}

\begin{tcolorbox}[breakable,colback=teal!4,colframe=teal!35!black,title={Before you move on}]
What does der Sarg mean? (\textquotedblleft{}The coffin\textquotedblright{}.) Say it once more.
\end{tcolorbox}
\section[der Friedhof]{der Friedhof — the cemetery, the graveyard}

\label{lesson:GE-C303-friedhof}



\begin{quote}
Before the new one: say the German for the funeral, then the German for the coffin.
\end{quote}

\subsection*{You'll want to know: der Friedhof}
\textbf{der Friedhof} — \textquotedblleft{}the cemetery, the graveyard\textquotedblright{}.

\begin{quote}
Der Friedhof liegt hinter der Kirche.
\end{quote}

\begin{quote}
\emph{The cemetery lies behind the church.}
\end{quote}

Plural: \emph{die Friedhöfe}.

\begin{cousinweb}
\emph{Friedhof} is inherited from Old High German \emph{frīthof} \textquotedblleft{}an enclosed yard\textquotedblright{} (EtymDB): the first part is usually traced to an old word for an enclosure, later heard as \emph{Frieden} \textquotedblleft{}peace\textquotedblright{}.
\end{cousinweb}

\begin{grammarlens}[title={What you've built}]
One more word for telling, explaining and describing at B1.
\end{grammarlens}

\subsection*{Your turn}
\begin{itemize}
\raggedright
  \item \emph{Say it:} \emph{der Friedhof}
  \item \emph{Say it:} \emph{der Friedhof}, once more
  \item \emph{Say it:} a sentence of your own with \emph{der Friedhof}
  \item \emph{Recall:} say the German for the funeral, then the German for the coffin, then say \emph{der Friedhof} again
\end{itemize}

\begin{tcolorbox}[breakable,colback=teal!4,colframe=teal!35!black,title={Before you move on}]
What does der Friedhof mean? (\textquotedblleft{}The cemetery, the graveyard\textquotedblright{}.) Say it once more.
\end{tcolorbox}
